% 第九卷 14.1.1.1：GTS目标驱动树解码在金额题上的生成过程（正文栏宽，约105 mm）。
% 表达式树 Add(Mul(Mul(25,3),0.8), Mul(4,5))，数量按原文顺序记为 n0=3, n1=25, n2=5, n3=4, n4=0.8。
% 节点左上角数字为前序生成的步次；蓝色实线：父目标向下分解出左、右子目标；
% 红色虚线：左子树完成后自底向上合成的子树嵌入 t，送去生成右兄弟的目标。
% FigureModel：运算符节点；FigureInput：数量叶节点。
\begin{tikzpicture}[
  x=1mm,y=1mm,
  figure picture,
  op/.style={figure node,draw=FigureModel,fill=FigureModelFill,minimum width=9mm,
    minimum height=6.5mm,inner sep=1pt},
  leaf/.style={figure node,draw=FigureInput,fill=FigureInputFill,minimum width=11mm,
    minimum height=6.5mm,inner sep=1pt},
  goal/.style={figure flow,draw=FigureInput},
  emb/.style={figure flow,draw=FigureOutput,dash pattern=on 2pt off 1.5pt},
  stepno/.style={figure note,inner sep=0.5pt}
]
  \node[op] (r) at (56,42) {$+$};
  \node[op] (a) at (28,29) {$\times$};
  \node[op] (b) at (84,29) {$\times$};
  \node[op] (c) at (14,16) {$\times$};
  \node[leaf] (n4) at (42,16) {$n_4{=}0.8$};
  \node[leaf] (n1) at (5,3) {$n_1{=}25$};
  \node[leaf] (n0) at (23,3) {$n_0{=}3$};
  \node[leaf] (n3) at (75,16) {$n_3{=}4$};
  \node[leaf] (n2) at (93,16) {$n_2{=}5$};

  \foreach \n/\k in {r/1,a/2,c/3,n1/4,n0/5,n4/6,b/7,n3/8,n2/9}{
    \node[stepno,anchor=south east] at (\n.north west) {\k};
  }

  \draw[goal] (r) -- node[figure edge label,above left=-1pt] {$\bm{q}_{\mathrm{l}}$} (a);
  \draw[goal] (r) -- node[figure edge label,above right=-1pt] {$\bm{q}_{\mathrm{r}}$} (b);
  \draw[goal] (a) -- (c);
  \draw[goal] (a) -- (n4);
  \draw[goal] (c) -- (n1);
  \draw[goal] (c) -- (n0);
  \draw[goal] (b) -- (n3);
  \draw[goal] (b) -- (n2);

  \draw[emb] (c.east) to[bend right=25] node[figure edge label,below=1pt] {$\bm{t}_{\mathrm{l}}$} (n4.west);
  \draw[emb] (a.east) to[bend right=12] node[figure edge label,below=1pt] {$\bm{t}_{\mathrm{l}}$} (b.west);

  \node[figure note,anchor=west,align=left] at (60,4)
    {前序：$+\ \times\ \times\ n_1\ n_0\ n_4\ \times\ n_3\ n_2$};
\end{tikzpicture}
